\section{Related Work}
\paragraph{Multi-Agent Systems}
Inspired by collaborative problem solving, early efforts propose multi-agent systems (MAS) to enhance the task-solving capability of language models~\cite{zhang2025agentorchestra, DBLP:journals/corr/abs-2505-07437,wu2024autogen, shi2025aime, gao2025single, zhu2025oagents, fang2025cognitive, zhang2024aflow,DBLP:conf/icml/Li0FXC0L25}.
For example, MetaGPT~\cite{hong2023metagpt} organizes agents into a structured software-development workflow, where specialized roles (e.g., product manager, architect) collaborate via predefined communication protocols.
OWL~\cite{hu2025owl} adopts a planner-worker workflow to improve transfer and generalization by modularizing domain-agnostic planning and domain-specific execution.
Despite their effectiveness, most MAS typically rely on a fixed workflow, leading to rigidity.
Although AutoAgents~\cite{chen2023autoagents} proposes building different multi-agent systems for each task, they still rely on a fixed workflow to accomplish this.
This motivates a growing shift toward the \emph{sub-agents-as-tools} paradigm, and we will list related works in the next part~\cite{gao2025evolvesurvey, gao2025single}. 
\our{} follows the latter and further emphasizes orchestration-centric, dynamic sub-agent creation without relying on a specific human-designed workflow. 

\paragraph{Sub-Agent as Tools}
This approach involves a primary agentic model invoking a sub-agent in a tool-like manner to solve problems~\cite{li2025chain, su2025toolorchestra, grand2025self,DBLP:journals/tkde/LiuSLMJZFLTL25}. For example, THREAD~\cite{schroeder2025thread} enables the recursive spawning of sub-agents to address decomposed subproblems. Similarly, Context-Folding~\cite{sun2025scaling} proposes branching for a subtask and then folding it back by compressing intermediate steps into a concise summary, thereby managing context.
However, these methods do not treat sub-agents as fully specialized agents, leading to their insufficient utilization. 
Other practical systems, such as Claude Code~\cite{anthropic_claude_code_subagents}, support sub-agents that operate within isolated context windows with custom system prompts and tool permissions. Yet, these sub-agents are typically configured as fixed specialists and still require manual design.
\our{} addresses these limitations by treating each sub-agent as a dynamic unit and proposes an orchestration-centric agentic system that proactively and dynamically creates such sub-agents on demand.

